% Chapter 3 — The Objects (KLEPPMANN). Budget 8.5 pp (7 objects + 1.5 Term Sheets as Graphs).
\section{The Objects}\label{ch:objects}

This chapter discharges \S 4 of the constitution as amended, including Amendment~2 (the
mass-never-value sentence) and Amendment~3 (the signed coordinate basis).

At the highest level a ledger is three things and a record: units --- what can be held;
wallets --- where it is held; balances --- how much of each unit each wallet holds; and
the immutable log that records every change to them. Each object is defined here once;
nothing later in the specification introduces another. The chapter ends where a unit's
terms begin: with the product graph, the declared structure from which a unit's watches,
its state schema, and its contract's actions are all derivable.

\subsection{Units}

A \textbf{unit} is anything that can be held or owed inside the ledger: a currency (USD
in cents), an equity security identified by its ISIN, a bond, a listed or OTC derivative,
a securities loan, a managed-account mandate, a margin obligation under a collateral
agreement, or any other contractual instrument. Units are first-class and extensible: a
new instrument type is a new element of the universe of units, never a modification of
the machinery. Because every unit --- asset, liability, or contractual obligation --- is
represented uniformly, one move machinery and one conservation discipline cover them all,
without special cases.

Units divide into \textbf{valued} and \textbf{non-valued}. A unit is valued exactly when
it represents an economic claim whose value exists independently of the governance or
continued operation of the instrument that created it: a stock, a bond, a client's
redemption claim on a managed account, a total return swap. A unit is non-valued when it
is the governance itself: a mandate's rulebook, a portfolio-swap reset mechanism, a
quantitative strategy's governance shell. The price of a non-valued unit is undefined ---
not zero --- because pricing the governance structure beside the claim it governs would
count the same economic reality twice. Valuation (Chapter~7) sums over the valued units
only.

Every unit carries three things: its \textbf{terms} --- the contractual text, versioned
and append-only; its \textbf{state} --- where it stands in its lifecycle; and its
\textbf{smart contracts} --- the executable form of those terms. From the terms follow
its \textbf{watches}: the dates, observations, and conditions the Event Monitor must
observe on its behalf. Declaring them is the unit's duty --- at registration, and again
whenever a later transaction creates a new obligation --- and the declaration is
exhaustive: the Monitor does exactly what the units have asked of it, nothing more and
nothing less. Section~\ref{sec:termsheet-graphs} shows that this duty is not a discipline
to remember but a consequence of the form the terms take.

\medskip\noindent\textbf{Episode --- the future (T1): terms declare the watches.}
\emph{Trigger.} The exchange lists FUT-IDX, and on day~0 W-ALPHA's purchase of one
contract at 995.00 executes; both facts are captured and recorded at the boundary.
\emph{Contract.} Each recorded event is mapped to its transaction by the responsible
contract (Chapter~2): the listing to a registration, the execution to a trade booking.
\emph{Transactions.} Two, in order:
\begin{enumerate}[nosep,leftmargin=2.2em]
\item \emph{Registration} (moveless). ProductTerms FUT-IDX version~1: cash-settled
  future on index IDX, multiplier 10, USD, settled-to-market --- the regime a declared
  term of the CCP agreement unit (Chapter~9). Unit-state change: watches declared ---
  one per daily settlement-price publication, and one for the expiry.
\item \emph{Trade} (day~0). One move: 1 FUT-IDX, owned, from the CCP's virtual wallet to
  W-ALPHA. One position-state change: W-ALPHA's traded price 995.00 recorded.
\end{enumerate}
\emph{The door.} Both transactions carry cause-derived identifiers, so a duplicated
capture commits once; the move's paired legs make conservation automatic; and no cash
moves --- a settled-to-market future exchanges nothing at trade, the first
variation-margin move being day~1's (Chapters~5 and~9).
\emph{Design point.} \emph{The watches follow from the terms and are declared at
registration: the future's whole life --- every print, the expiry --- is asked for
before the first event arrives.}

\medskip\noindent\textbf{Episode --- the variance swap (T5): the declaration is
exhaustive.}
\emph{Trigger.} VS-1 is executed and its confirmation captured at the boundary: a
one-year OTC variance swap on IDX, variance strike 400 variance points, variance
notional \USD{1{,}000} per point, no cap, 252 scheduled fixings against the recorded IDX
official closes --- the same observation source FUT-IDX settles against.
\emph{Contract.} The confirmation maps to the registration transaction that installs the
unit.
\emph{Transactions.} One, moveless:
\begin{enumerate}[nosep,leftmargin=2.2em]
\item ProductTerms VS-1 version~1: the terms above, the fixing schedule carried as
  declared data.
\item Unit-state change: 252 fixing watches declared --- one per scheduled fixing date,
  each a watch on the recorded IDX official close; the 252nd is the final fixing, on
  whose firing settlement is computed (Chapter~7).
\end{enumerate}
\emph{The door.} A moveless transaction is admitted like any other: idempotent under its
cause-derived identifier; nothing to conserve, everything to record.
\emph{Design point.} \emph{Exhaustive means exhaustive: all 252 conditions the swap's
life requires are on the record before the first fixing prints, and no watch is
discovered mid-life.}

\medskip\noindent\textbf{Episode --- the TRS (T6): a non-valued unit beside an ordinary
one.}
\emph{Trigger.} Two registrations are captured at the boundary: the strategy S-QIS, and
the swap TRS-1 written on the strategy's wallet.
\emph{Contract.} Each confirmation maps to its registration transaction.
\emph{Transactions.} Two, both moveless:
\begin{enumerate}[nosep,leftmargin=2.2em]
\item ProductTerms S-QIS version~1: the strategy's rulebook, the complete terms of a
  \emph{non-valued} unit whose contract manages wallet W-QIS in virtual ledger V-1
  (Chapter~10). S-QIS has no price: undefined, not zero.
\item ProductTerms TRS-1 version~1: notional \USD{10{,}000{,}000}; reference leg the NAV
  of W-QIS; financing leg fixed 4\% per annum, quarterly. Unit-state change: watches
  declared --- each quarterly reset date, at which the contract reads the stamped NAV
  observation then on the record.
\end{enumerate}
\emph{The door.} Consistency of reference: TRS-1's terms name a wallet and an
observation source that exist on the record; the registrations commit once under their
cause-derived identifiers.
\emph{Design point.} \emph{The strategy's governance shell is non-valued because pricing
it beside the wallet it governs would count the same economics twice; the TRS that
references it is an ordinary valued unit of the true ledger --- referencing a virtual
ledger makes nothing special.}

\subsection{Wallets}

A \textbf{wallet} is a named logical partition of position space --- a container, and
deliberately nothing more. Its state at any moment is, for each unit it holds, a signed
quantity: its \textbf{balance}. A wallet names two parties: a \textbf{manager},
authorised to perform transactions on it, and a \textbf{beneficial owner}, legally
accountable for the book. Both are declared at creation, and the distinction is
load-bearing: authorisation is checked against the manager, accountability and reporting
attach to the beneficial owner. Wallets carry no smart contract and no logic of their
own; where a wallet appears to need behaviour --- the rules of a managed account ---
that behaviour is a unit (the mandate), and the wallet stays a container. W-QIS above is
exactly this shape: the rulebook lives in the non-valued unit S-QIS, and the wallet only
holds.

\subsection{The Atomic Move}

An \textbf{atomic move} is the sole operation that changes a balance: it transfers a
positive quantity of exactly one unit from a source wallet to a destination wallet in a
single indivisible step --- the same quantity debited at the source and credited at the
destination. Because a move only transfers, it can never create or destroy quantity;
issuance and retirement are themselves paired-leg moves against the issuer's wallet
(Chapter~14 states the conservation law this yields and proves it by induction on the
log).

Quantities are exact integers in minor units: a USD amount is an integer count of cents,
a holding an integer count of shares or contracts, a recorded statistic an integer at
its declared scale. No quantity in the ledger is a floating-point number. Where an event
divides unevenly --- a per-share figure across a holding it does not divide, a pro-rata
allocation with an odd cent --- the rounding convention is part of the declared terms,
and the remainder is booked explicitly as its own move: the convention says where the
odd cent goes, and the odd cent's journey is on the record. Nothing is ever silently
dropped, which is what makes conservation exact rather than approximate.

\subsection{The Transaction}

A \textbf{transaction} is an atomic, all-or-nothing bundle of four kinds of change:
moves, unit-state changes, position-state changes, and terms changes. The four align
one-for-one with the four components of the ledger's state: balances, and the three
homes of Chapter~6 --- UnitStatus, PositionState, ProductTerms. A transaction applies in
full or not at all, and some transactions carry no moves: the registration of a new unit
and the recording of a barrier breach are complete transactions with empty move lists.
The episodes above have already shown three of the four kinds; the split shows the
fourth.

\medskip\noindent\textbf{Episode --- the split (T4): a terms change is a transaction
component.}
\emph{Trigger.} 2026-06-01: ACME's 2-for-1 split becomes effective. The recorded split
notice had declared the effective-date watch; the condition is met and the Event Monitor
emits the event.
\emph{Contract.} ACME's contract fires with the event and the current projection of the
ledger: W-ALPHA holds 10{,}000 shares.
\emph{Transactions.} One, atomic:
\begin{enumerate}[nosep,leftmargin=2.2em]
\item Move: 10{,}000 ACME, owned, from the issuer's virtual wallet to W-ALPHA --- the
  holding becomes 20{,}000 by a paired-leg move, one such move per holder, never by an
  edit.
\item Terms change: ProductTerms ACME version~2, carrying the 2-for-1 event. How the
  recorded price frame 100.00 reads as 50.00 after the event --- and the declared
  per-share dividend figure 1.50 as 0.75 --- is the market data operator's work,
  Chapter~8.
\end{enumerate}
\emph{The door.} Atomicity: the move and the new terms version apply in full or not at
all --- no instant exists at which 20{,}000 shares stand beside pre-split terms
(Chapter~14's consistency of reference begins here). Conservation: the issuer wallet's
balance falls by 10{,}000 as W-ALPHA's rises.
\emph{Design point.} \emph{A terms change is one of the four components of a
transaction, atomic with the moves it accompanies; terms are versioned and append-only,
never edited.}

\subsection{The Immutable Log}

The \textbf{immutable log} is the ordered, append-only, hash-chained record of every
admitted transaction --- hash-chained so that any alteration of history is detectable.
It is the sole canonical source of truth. Three words are used precisely throughout this
specification. The \textbf{log} is the recorded sequence itself. The \textbf{ledger} is
the state obtained by folding the log: balances and the three homes. The \textbf{record}
is everything that has been recorded --- the log and all it carries: transactions,
events, observations. ``On the record'' means present there, and a \textbf{projection}
is any view computed from the record --- which is why, in Chapter~2's typed picture, the
watch arrow reads the Record while contract and apply consume the Ledger.

\subsection{Closure by Virtual Wallets}

Every counterparty outside the system is represented by a \textbf{virtual wallet} inside
it: the CCP a future faces, the issuer of a stock, the writer of an option each hold a
wallet, and every move has both legs on the book. The system is therefore closed, and
conservation holds by construction, not by check. The split episode shows the closure
working: the 10{,}000 new shares credited to W-ALPHA are debited from the issuer's
virtual wallet in the same move --- quantity appears nowhere from nothing, and the sum
over all wallets, virtual ones included, stays at zero.

\subsection{The Coordinate Vector}

When a balance must carry more information than a single number, it generalises to a
signed vector of coordinates: \textbf{(owned, lent, posted)}. The five constitutional
names are the rays of this basis --- \emph{lent out} and \emph{borrowed} are the two
signs of lent; \emph{posted as collateral} and \emph{received as collateral} are the two
signs of posted --- and the outright holding is the degenerate case with all mass on
owned. A quantity earns a coordinate only when a distinct real-world action can change
that coordinate independently of ownership: lending changes possession without changing
ownership, pledging encumbers without transferring, and nothing else in scope passes the
test, so there is no fourth coordinate. A move writes exactly one coordinate of one
unit, paired legs as always, so conservation holds per (unit, coordinate), not merely
per unit.

Only the owned coordinate carries economic value and drives profit and loss. The
remaining coordinates carry mass, never value: they record possession and encumbrance
under a named agreement, no quantity reaches them without reference to the agreement
unit that governs it, and a lifecycle event extinguishes value, never mass --- a unit
leaves the ledger only from the zero vector. Anything computable from the coordinates
--- availability, encumbrance, possession --- is a projection, computed when needed and
never stored.

This section fixes the definition and nothing else. Which coordinate a given inflow
writes, the agreement semantics of the non-owned planes, the guards checked at the door,
and the regime key that decides among them are Chapter~9's.

\subsection{Term Sheets as Graphs}\label{sec:termsheet-graphs}

A unit's terms are not free text with a contract bolted on; they have a shape, and the
shape is a graph. Take the one-touch OT-1 of the thread map: payout \USD{1{,}000{,}000},
its barrier breached when OMEGA's official close is at or below 80.00 --- a
discretely-monitored (closing-price) one-touch. Its life has three states --- live,
triggered, expired. From live, exactly two things can happen: a recorded OMEGA official
close at or below 80.00 (the unit becomes triggered and the payout falls due), or the
expiry date arrives untouched (the unit becomes expired, worthless). From triggered and
from expired, nothing further happens. Drawn, that is a graph: three nodes, two edges
leaving one of them.

\begin{definition}[Product graph]\label{def:product-graph}
The terms of a unit define its \emph{product graph}: the nodes are the unit's lifecycle
states, a set closed and known at registration; each node carries a payload --- the
sufficiency facts of constitution \S 7, held in the three homes of Chapter~6 (the
variance swap's accrued fixing statistic, the future's last settlement applied); each
edge carries a \emph{guard} --- an event kind the Event Monitor can watch, together with
a predicate --- and an \emph{action} --- the transaction template the traversal emits.
\end{definition}

The structure is declared data: nodes, edges, guards, and payload schemas are recorded
with the terms and readable without running anything. The arithmetic inside an action is
not data and does not pretend to be: it is a small, named, versioned pure function ---
the variance swap's realised-variance accumulator, the future's margin difference ---
referenced from the edge that uses it. Declared structure carries behaviour, per the
constitution's Clarity commitment (\S 2); named pure functions carry computation; and
nothing is hidden in code the graph does not point to.

The graph is identical with the watch discipline of this chapter. The watch list of a
unit is exactly the out-edge set of its current node, so declaring the graph declares
the watches --- the exhaustiveness demanded of units above is not a rule to remember but
a projection of the graph. Arming is edge registration; firing is traversal; and a
traversal expires its sibling edges: when OT-1's touch edge fires, the expiry edge dies
with the live node that carried it. The declared~$\to$~armed~$\to$~fired/expired
lifecycle of Chapter~4's handshake is a graph fact, not a bookkeeping convention. The
variance swap reads the same way: 252 fixing edges, and the 252 declared watches of its
episode are their guards, exhaustive at registration because the graph is closed at
registration.

Because the node set is closed at registration, each state can be a distinct type and
each edge a function between those types --- and the edges are the only constructors
those types have:

\begin{lstlisting}
data Live      = Live      { accrued  :: Payload }
data Triggered = Triggered { touchRef :: EventId }
data Expired   = Expired

touch  :: Fixing -> Live -> Triggered   -- the only constructor of Triggered
expire :: Expiry -> Live -> Expired
-- Triggered -> Live: no such function exists, so no program un-knocks a
-- note; a case analysis that forgets a node does not compile.
\end{lstlisting}

An illegal transition is not refused; it is unrepresentable --- there is no program that
writes it. Totality is compiler-checked: exhaustive case analysis over a closed node set
is exactly what these types force, so a contract that forgets the expired state is a
compile error, not a production incident.

One requirement makes the identification normative rather than decorative.

\begin{invariant}[Graph consistency]\label{inv:graph-consistency}
For every unit there exists a product graph from which its watches, its state schema,
and its contract's actions are all derivable, and their mutual consistency holds:
(i)~at every commit, the unit's watch list equals the out-edge set of its current node;
(ii)~every fired event matches exactly one out-edge and lands the unit on that edge's
target node;
(iii)~every transaction the contract emits equals the edge's declared action applied to
the recorded inputs;
(iv)~a traversal expires exactly the declared sibling edges, and no others.
\end{invariant}

The four properties are executable, not defended in prose: their executable forms join
the invariant catalogue of Chapter~14 and the test regime of Chapter~15, where they are
property-tested over generated products, events, and histories.

\begin{secondtelling}
The product graph is one object, presented by its generators (nodes and edges) and the
relations the invariant states. The Event Monitor's watch list, the typed state schema,
and the contract's action table are three interpreters of that object ---
structure-preserving maps out of the same graph --- and the consistency invariant says
exactly that the three maps agree on it. Nothing interprets a term sheet except through
its graph.
\end{secondtelling}
